\chapter{Équations aux dérivées partielles}

Les équations aux dérivées partielles décrivent de nombreux problèmes de
physique et de mécanique. Leur résolution numérique fait appel à des formules
de discrétisation qui convertissent l'équation aux dérivées partielles en un
système d'équations discrètes que l'on peut résoudre par des schémas
itératifs.

Ce chapitre présente les méthodes de base de discrétisation des équations aux
dérivées partielles. En particulier, on s'intéresse ici aux équations de
Poisson et de Laplace, à l'équation de la chaleur et à l'équation d'ondes.
Pour que les schémas discrets conduisent à des solutions viables, ils doivent
être stables, c'est-à-dire que les solutions numériques doivent être très peu
sensibles au changement des pas de discrétisation. C'est pourquoi, ce chapitre
s'intéresse aussi à l'établissement des conditions de stabilité des schémas
discrets.

\section{Notions générales}

\subsection{Classification}

Considérons une équation aux dérivées partielles d'ordre 2 de la forme
\[
  A\frac{\partial^2 U}{\partial x^2} + 2B\frac{\partial^2 U}{\partial x\,\partial y}
  + C\frac{\partial^2 U}{\partial y^2}
  + F\Bigl(x, y, U, \frac{\partial U}{\partial x}, \frac{\partial U}{\partial y}\Bigr) = 0
\]
où $A$, $B$ et $C$ sont des fonctions de $x$ et $y$ et sont de classe $C^2$.
\begin{itemize}[label=$\blacklozenge$]
  \item Si $B^2 - AC = 0 \;\Longrightarrow\;$ équation \emph{parabolique} ;
  \item Si $B^2 - AC < 0 \;\Longrightarrow\;$ équation \emph{elliptique} ;
  \item Si $B^2 - AC > 0 \;\Longrightarrow\;$ équation \emph{hyperbolique}.
\end{itemize}

Pour résoudre une équation aux dérivées partielles, il faut connaître les
conditions aux limites ou aux frontières et les conditions initiales. La
nature de ces conditions permet de définir les types de problèmes suivants :

\textbf{Problème de DIRICHLET pur :} les valeurs de $U$ sur les frontières du
domaine sont connues.

\textbf{Problème de NEUMANN pur :} les valeurs des dérivées de $U$ sur les
frontières sont connues.

\textbf{Problème mixte DIRICHLET-NEUMANN :} les valeurs de $U$ sont connues
sur une partie des frontières et celles des dérivées de $U$ sur la partie
restante.

\textbf{Problème de CAUCHY :} les valeurs de $U$ et de ses dérivées sont
connues sur les frontières du domaine.

\subsection{Quelques équations aux dérivées partielles en Physique}

\subsubsection{Équation d'onde}

Elle a la forme :
\[
  \frac{\partial^2 U}{\partial t^2} = a^2\Delta U
  + f\Bigl(\frac{\partial U}{\partial t}, x, y, z, t\Bigr)
\]
où $a$ est la vitesse de l'onde et $f$ une action extérieure.

\textbf{En dimension 1}, cette équation décrit :
\begin{itemize}[label=--]
  \item Les vibrations transversales d'une corde : l'équation générale dans ce
        cas est
        \[
          \rho(x)\frac{\partial^2 U}{\partial t^2} + k\frac{\partial U}{\partial t}
          = T\frac{\partial^2 U}{\partial x^2} + f(x,t)
        \]
        où $\rho$ est la masse volumique, $k$ le coefficient de dissipation,
        $T$ la tension dans la corde. Ici la vitesse des ondes est
        $a = \sqrt{T/\rho}$.
  \item Les vibrations longitudinales dans une barre : $a = \sqrt{E/\rho}$ où
        $E$ est le module d'élasticité.
  \item Les vibrations de torsion d'une barre avec $a = \sqrt{\mu/\rho}$ où
        $\mu = \dfrac{E}{2(1+\sigma)}$ et $\sigma$ est le coefficient de
        Poisson.
  \item La propagation d'ondes électromagnétiques ou électriques.
\end{itemize}

\textbf{En dimension 2}, elle peut décrire les vibrations d'une membrane
plane.

\textbf{En dimension 3}, elle décrit les ondes sonores en hydrodynamique, les
ondes électromagnétiques (cas des cavités résonantes et du rayonnement des
antennes).

\subsubsection{Équation de Laplace et de Poisson}

Elles s'écrivent sous la forme suivante
\[
  \begin{aligned}
    \Delta U &= 0 && \text{(Laplace)}\\
    \Delta U &= f(x, y, z) && \text{(Poisson)}
  \end{aligned}
\]
Ce sont les prototypes des équations elliptiques. Elles décrivent les
problèmes stationnaires suivants :
\begin{itemize}[label=--]
  \item répartition des charges électriques sur la surface d'un conducteur ;
  \item équilibre thermique dans un corps homogène ;
  \item distribution du potentiel des vitesses dans un fluide non visqueux,
        incompressible et en écoulement irrotationnel.
\end{itemize}

\subsubsection{Équation de la chaleur}

Sa forme générale est :
\[
  c\rho\frac{\partial U}{\partial t} = \operatorname{div}\bigl(k\,\overrightarrow{\operatorname{grad}}\,U\bigr)
  + F(x, y, z, t)
\]
où $U$ est la température en un point $P(x, y, z)$ au temps $t$, $\rho$ la
masse volumique, $c$ la chaleur spécifique, $k$ la conductivité interne et $F$
l'apport des sources ou des puits de chaleur.

Si $c$, $\rho$ et $k$ sont constantes, cette équation se réduit à :
\[
  \frac{\partial U}{\partial t} = a^2\Delta U + f(x, y, z, t) .
\]
Dans le cas où
\[
  \frac{\partial U}{\partial t} = 0,
\]
on est en présence d'un problème de transfert de chaleur en régime permanent
(conduction et convection).

Dans le cas général où $\dfrac{\partial U}{\partial t} \neq 0$, on est en
présence d'un problème de diffusion de la chaleur.

\section{Équation de Poisson dans le plan}

\subsection{Discrétisation de l'équation de Poisson}

Soit l'équation :
\begin{equation}
  \frac{\partial^2 U}{\partial x^2} + \frac{\partial^2 U}{\partial y^2} = f(x, y)
  \tag{P}
\end{equation}
définie sur un domaine $D$ rectangulaire :
$D = \{(x, y) \in \R^2 \,/\, a < x < b \text{ et } c < y < d\}$.

Soit $(S)$ la frontière de $D$, la condition suivante :
\[
  U(x, y) = g(x, y)
\]
est imposée sur $(S)$.

Nous considérons que les fonctions $f$ et $g$ sont continues sur le domaine
$D$. Soient $h$ et $k$ les pas de calcul suivant $x$ et $y$ respectivement.
Soient $n$ et $m$ les nombres entiers définis par
\[
  h = \frac{b - a}{n} \qquad\text{et}\qquad k = \frac{d - c}{m} .
\]
Un point du domaine $D$ est défini par les coordonnées suivantes :
\[
  \left\{
  \begin{aligned}
    x_i &= a + ih && \text{avec } i = 1, \dots, n-1\\
    y_j &= c + jk && \text{avec } j = 1, \dots, m-1
  \end{aligned}
  \right.
\]
D'après les définitions des dérivées, on a :
\begin{equation}
  \frac{\partial^2 U}{\partial x^2}
  = \frac{U_{i+1,j} - 2U_{i,j} + U_{i-1,j}}{h^2}
  - \frac{h^2}{12}\frac{\partial^4 U}{\partial x^4}(\xi_{ij}, y_j)
  \tag{P1}
\end{equation}
où $U_{ij} = U(x_i, y_j)$ ; $\xi \in\, ]x_{i-1}, x_{i+1}[$.

En écrivant une équation similaire pour $\dfrac{\partial^2 U}{\partial y^2}$,
l'équation (P) devient, avec une erreur de troncature d'ordre
$O(h^2 + k^2)$, le système suivant :
\begin{equation}
  2\left[\Bigl(\frac{h}{k}\Bigr)^2 + 1\right]U_{i,j} - \bigl(U_{i+1,j} + U_{i-1,j}\bigr)
  - \Bigl(\frac{h}{k}\Bigr)^2\bigl(U_{i,j+1} + U_{i,j-1}\bigr) = -h^2f_{i,j}
  \tag{P2}
\end{equation}
où $f_{i,j} = f(x_i, y_j)$ ; avec $i = 1, 2, \dots, n-1$ et
$j = 1, 2, \dots, m-1$.

La discrétisation des conditions aux frontières est la suivante :
\begin{equation}
  \left\{
  \begin{aligned}
    U_{0,j} &= g(x_0, y_j)\\
    U_{n,j} &= g(x_n, y_j)
  \end{aligned}
  \right.
  \quad j = 0, 1, \dots, m
  \qquad\qquad
  \left\{
  \begin{aligned}
    U_{i,0} &= g(x_i, y_0)\\
    U_{i,m} &= g(x_i, y_m)
  \end{aligned}
  \right.
  \quad i = 0, 1, \dots, n
  \tag{P3}
\end{equation}
En associant les équations (P2) et (P3), on obtient le système d'équations
algébriques linéaires en $U_{ij}$ que l'on peut résoudre par la méthode
d'élimination de Gauss.

\begin{exercice}
Écrire le schéma d'itération de Gauss-Seidel et de Jacobi pour résoudre
l'équation de Poisson.
\end{exercice}

\subsection{Cas spécifiques}

\subsubsection{Conditions aux limites de type Neumann}

Lorsque les conditions aux limites sont de type Neumann, le schéma de
discrétisation doit se faire avec beaucoup de précautions. Généralement, il
est conseillé de combiner cette condition avec l'équation aux dérivées
partielles de départ.

Par exemple, si la condition de Neumann s'écrit
$\left.\dfrac{\partial U}{\partial y}\right|_S = g(x, y)$, alors on utilise le
schéma suivant. On développe $U(x_i, y_1)$ autour du point $(x_i, y_0)$ ;
c'est-à-dire que l'on écrit :
\[
  U(x_i, y_1) = U(x_i, y_0) + k\frac{\partial U}{\partial y}(x_i, y_0)
  + \frac{k^2}{2}\frac{\partial^2 U}{\partial y^2}(x_i, y_0) + O(k^3) .
\]
La quantité $\dfrac{\partial^2 U}{\partial y^2}$ en $(x_i, y_0)$ peut être
remplacée par son équivalent à partir de l'équation aux dérivées partielles
de départ. Par exemple, dans le cas de l'équation de Laplace, on a :
$\dfrac{\partial^2 U}{\partial y^2} = -\dfrac{\partial^2 U}{\partial x^2}$.
Il vient alors :
\[
  \begin{aligned}
    U(x_i, y_1) &= U(x_i, y_0) + k\frac{\partial U}{\partial y}(x_i, y_0)
                 - \frac{k^2}{2}\frac{\partial^2 U}{\partial x^2}(x_i, y_0)\\
                &= U(x_i, y_0) + k\frac{\partial U}{\partial y}(x_i, y_0)
                 - \frac{k^2}{2h^2}\bigl(U_{i+1,0} - 2U_{i,0} + U_{i-1,0}\bigr).
  \end{aligned}
\]
Finalement la condition de Neumann donne :
\begin{equation}
  U_{i,1} - \Bigl(1 + \frac{k^2}{h^2}\Bigr)U_{i,0}
  + \frac{k^2}{2h^2}\bigl(U_{i+1,0} + U_{i-1,0}\bigr) = kg(x_i, y_0) .
  \tag{P4}
\end{equation}
On peut alors associer à (P4) l'équation (P2) pour avoir le système discret
global à résoudre.

\subsubsection{Théorème du maximum et du minimum discrets}

Dans le cas de l'équation de Laplace, on a :
\[
  U_{i,j} = \frac{k^2}{2(k^2 + h^2)}\bigl(U_{i+1,j} + U_{i-1,j}\bigr)
          + \frac{h^2}{2(k^2 + h^2)}\bigl(U_{i,j+1} + U_{i,j-1}\bigr).
\]
Donc $U_{ij}$ est le barycentre des nombres $U_{i+1,j}$, $U_{i-1,j}$,
$U_{i,j+1}$ et $U_{i,j-1}$ affectés des poids respectifs :
$\alpha = \dfrac{k^2}{2(k^2 + h^2)}$, $\alpha$,
$\beta = \dfrac{h^2}{2(k^2 + h^2)}$ et $\beta$. Puisque ces poids sont tous
positifs, $U_{ij}$ est donc compris entre la plus grande et la plus petite des
quantités $U_{i+1,j}$, $U_{i-1,j}$, $U_{i,j+1}$ et $U_{i,j-1}$. D'où le
théorème suivant :

\paragraph{\underline{Théorème}} Dans le cas de l'équation de Laplace, il n'y
a pas de maximum et de minimum stricts de $U$ à l'intérieur du domaine $D$. Le
maximum et le minimum sont atteints sur les frontières de $D$.

\paragraph{\underline{Remarque}} D'autres schémas de différences finies
peuvent être établis à partir des formules de discrétisation des dérivées. Par
exemple, on peut définir $\dfrac{\partial^2 U}{\partial x^2}$ par des schémas
présentant des erreurs d'ordre $O(h^4 + k^4)$.

\section{Équation de la chaleur}

\subsection{Problème à une dimension d'espace}

Nous nous limiterons au cas unidimensionnel. On a l'équation :
\begin{equation}
  \frac{\partial U}{\partial t} = a^2\frac{\partial^2 U}{\partial x^2}
  \tag{C1}
\end{equation}
dans le domaine $D = T \times L$ ; où $T = \R_+$ et $L = \,]0, l[$.

Les conditions aux limites sont : $U(0, t) = 0$, $U(l, t) = 0$.

Et la condition initiale est
\[
  U(x, 0) = f(x) .
\]
Il s'agit par exemple d'une tige de longueur $l$ dont les extrémités $x = 0$
et $x = l$ sont maintenues à la température $U = 0$. À l'instant initial, on
la soumet à une distribution de température $U(x, 0) = f(x)$. Il est question
de suivre l'évolution temporelle de la température dans la poutre.

\subsection{Méthode des différences progressives ou schéma explicite}

Soit $h$ le pas spatial et $k$ le pas temporel et soit $m$ le nombre défini
par $mh = l$.

On a
\[
  \frac{\partial U}{\partial t} = \frac{U_{i,j+1} - U_{i,j}}{k}
  \qquad\text{et}\qquad
  \frac{\partial^2 U}{\partial x^2} = \frac{U_{i+1,j} - 2U_{i,j} + U_{i-1,j}}{h^2} .
\]
Dans le domaine $D$, on obtient alors le système discret suivant :
\begin{equation}
  U_{i,j+1} = \Bigl(1 - \frac{2a^2k}{h^2}\Bigr)U_{i,j}
            + \frac{a^2k}{h^2}\bigl(U_{i+1,j} + U_{i-1,j}\bigr) ;
  \quad i = 1, 2, \dots, m-1 ; \quad j = 0, 1, 2, \dots
  \tag{C2}
\end{equation}
Le schéma (C2) présente une erreur de troncature d'ordre $O(k + h^2)$.

C'est un schéma explicite car pour calculer $U_{i,j+1}$, il suffit de
remplacer les quantités figurant au second membre par leurs valeurs connues et
stockées en mémoire.

Les conditions initiales et aux limites se discrétisent de la manière
suivante :
\begin{equation}
  U_{i,0} = f(x_i), \; i = 0, 1, \dots, m \,; \qquad
  U_{0,j} = 0 \;\text{ et }\; U_{m,j} = 0, \; j = 0, 1, 2, \dots
  \tag{C3}
\end{equation}
La première condition (C3) peut être substituée dans (C2) pour déterminer
$U_{i,1}$ pour $i$ allant de $1$ jusqu'à $m-1$. La condition $U_{0,j} = 0$ et
$U_{m,j} = 0 \Longrightarrow U_{0,1} = U_{m,1} = 0$. On a alors toutes les
valeurs de $U_{i,1}$. En procédant de la même manière, on obtient les valeurs
de $U_{i,2}$, $U_{i,3}$ et ainsi de suite.

L'inconvénient de cette méthode est qu'elle est conditionnellement stable. Il
faut donc avoir une contrainte entre les pas $h$ et $k$ et les coefficients de
l'équation pour obtenir des résultats satisfaisants. C'est la condition de
stabilité que nous établissons dans le paragraphe qui suit.

\subsection{Condition de stabilité du schéma explicite}

Considérons le schéma aux différences finies progressives de l'équation de la
chaleur. Soit une perturbation $\varepsilon_{ij}$ de la solution numérique
$U_{i,j}$. La solution numérique devient alors
$U'_{i,j} = U_{i,j} + \varepsilon_{ij}$. En écrivant que $U'_{ij}$ vérifie le
schéma discret, on obtient que la perturbation $\varepsilon_{ij}$ vérifie
aussi l'équation discrète :
\[
  \frac{\varepsilon_{i,j+1} - \varepsilon_{i,j}}{k}
  = a^2\frac{\varepsilon_{i+1,j} - 2\varepsilon_{i,j} + \varepsilon_{i-1,j}}{h^2} .
\]
Le schéma discret est stable si $\varepsilon_{ij}$ reste borné ou tend vers 0
au fur et à mesure que le temps avance. Pour avoir la condition de stabilité,
nous supposons que
\[
  \varepsilon_{i,j} = \sin(iph)\,e^{-rjk} .
\]
Cette expression est utilisée en tenant compte de sa dépendance suivant
l'espace $x = ih$ et doit satisfaire aux conditions aux limites. En
introduisant cette expression dans l'équation vérifiée par $\varepsilon_{ij}$,
on obtient :
\[
  e^{-rk} = 1 - \frac{4ka^2}{h^2}\sin^2\frac{ph}{2} .
\]
Or $\varepsilon_{i,j+1} = \varepsilon_{i,j}\,e^{-rk}$.

Donc $e^{-rk}$ est le \emph{facteur d'amplification}.

Le schéma est stable si le facteur d'amplification est inférieur à 1 (en
module). Dans le cas particulier analysé ici, cette condition implique que :
\[
  \abs{1 - \frac{4ka^2}{h^2}\sin^2\frac{ph}{2}} < 1
  \;\Longrightarrow\;
  \abs{1 - \frac{4ka^2}{h^2}} < 1
  \;\Longrightarrow\;
  \frac{ka^2}{h^2} < \frac{1}{2} .
\]

\subsection{Autres schémas}

\subsubsection{Différences finies régressives ou schéma implicite}

Afin d'éviter l'utilisation d'une condition de stabilité, on peut utiliser les
différences finies régressives pour la dérivée temporelle ; c'est-à-dire
\begin{equation}
  \frac{\partial U}{\partial t} = \frac{U_{i,j} - U_{i,j-1}}{k} .
  \tag{C4}
\end{equation}
L'équation de la chaleur (C1) donne alors le système discret suivant :
\begin{equation}
  (1 + 2\lambda)\,U_{i,j} - \lambda U_{i+1,j} - \lambda U_{i-1,j} = U_{i,j-1}
  \qquad\text{pour}\quad 1 \le i \le m-1 \quad\text{et}\quad 1 \le j \le \infty
  \tag{C5}
\end{equation}
où $\lambda = a^2k/h^2$.

Compte tenu des conditions initiales et des conditions aux limites, on obtient
un système matriciel tri-diagonal que l'on peut résoudre par les méthodes
itératives ou la méthode d'élimination de Gauss.

Pour déterminer la condition de stabilité, on procède de la même façon que
précédemment. On trouve pour ce schéma implicite, que le facteur
d'amplification est
\[
  e^{-rk} = \frac{1}{1 + \dfrac{4ka^2}{h^2}\sin^2\dfrac{ph}{2}} .
\]
On note qu'il est toujours inférieur à 1. Par conséquent, ce schéma est
inconditionnellement stable.

\subsubsection{Méthode de Richardson}

Les méthodes ci-dessus ont une erreur de troncature d'ordre $O(k)$ en $k$. Si
on veut obtenir une erreur d'ordre $O(k^2)$, on utilise la formule suivante :
\begin{equation}
  \frac{\partial U}{\partial t} = \frac{U_{i,j+1} - U_{i,j-1}}{2k} .
  \tag{C6}
\end{equation}
Avec cette discrétisation de la dérivée temporelle, on aboutit à un schéma
instable. Il n'est donc pas utilisé dans le cadre de l'équation de la
chaleur.

\subsubsection{Méthode de Crank-Nicolson}

Cette méthode fait la somme du schéma des différences finies progressives et
du schéma des différences régressives. Mais il faut au préalable écrire la
formule des différences régressives à l'instant $j+1$. On obtient alors
\[
  \frac{U_{i,j+1} - U_{i,j}}{k}
  - \frac{a^2}{2h^2}\bigl(U_{i+1,j} - 2U_{i,j} + U_{i-1,j}
                         + U_{i+1,j+1} - 2U_{i,j+1} + U_{i-1,j+1}\bigr) = 0 .
\]
C'est une méthode plus précise qui présente une erreur d'ordre
$O(h^2 + k^2)$.

\subsubsection{Méthode de Dufort-Frankel}

Pour résoudre le problème d'instabilité de la méthode de Richardson, la
quantité $U_{ij}$ est remplacée par $\bigl(U_{i,j+1} + U_{i,j-1}\bigr)/2$ et
on obtient
\[
  \frac{U_{i,j+1} - U_{i,j-1}}{2k}
  - a^2\bigl(U_{i+1,j} - U_{i,j+1} - U_{i,j-1} + U_{i-1,j}\bigr)/h^2 = 0 .
\]
Ce schéma est inconditionnellement stable.

\begin{exercice}[Exercices]
\begin{enumerate}[label=\arabic*)]
  \item Établir l'algorithme pour le schéma explicite.
  \item Établir les conditions de stabilité des schémas de Crank-Nicolson et
        de Dufort-Frankel.
  \item Établir les algorithmes pour les méthodes des différences
        régressives, de Crank-Nicolson et de Dufort-Frankel.
\end{enumerate}
\end{exercice}

\section{Équation d'ondes}

Nous considérons le problème suivant :
\begin{equation}
  \frac{\partial^2 U}{\partial t^2} - a^2\frac{\partial^2 U}{\partial x^2} = 0
  \tag{O1}
\end{equation}
dans le domaine $D = T \times L$ et soumis aux conditions suivantes :
\[
  U(0, t) = U(l, t) = 0 \,; \qquad U(x, 0) = f(x) \,; \qquad
  \frac{\partial U(x, 0)}{\partial t} = g(x) .
\]
Il s'agit par exemple d'une corde vibrante fixée à ses deux extrémités et
excitée en une région par un déplacement transversal $U(x, 0) = f(x)$ et une
vitesse de déplacement transversal $\dfrac{\partial U(x, 0)}{\partial t} = g(x)$.
La question est de savoir comment cette perturbation va se propager le long de
la corde.

\subsection{Discrétisation directe}

La discrétisation de (O1) par les différences finies donne :
\begin{equation}
  U_{i,j+1} = 2(1 - \lambda^2)\,U_{i,j} + \lambda^2\bigl(U_{i+1,j} + U_{i-1,j}\bigr) - U_{i,j-1}
  \tag{O2}
\end{equation}
où $\lambda = ak/h$, $m = l/h$ et $i = 1, \dots, m-1$ ; $j = 1, \dots, \infty$.

Les conditions aux limites donnent :
\begin{equation}
  U_{0,j} = U_{m,j} = 0 .
  \tag{O3}
\end{equation}
La première condition initiale donne
\begin{equation}
  U_{i,0} = f(x_i) .
  \tag{O4}
\end{equation}
La deuxième condition initiale permet d'écrire :
\begin{equation}
  U_{i,1} = U_{i,0} + k\,g(x_i) .
  \tag{O5}
\end{equation}
Mais (O5) présente une erreur de troncature d'ordre $O(k)$. Pour avoir une
erreur d'ordre $O(k^2)$ comme l'équation (O2), on procède de la manière
suivante.
\[
  \frac{\partial U(x_i, 0)}{\partial t}
  = \frac{U(x_i, t_1) - U(x_i, 0)}{k}
  - \frac{k}{2}\frac{\partial^2 U(x_i, 0)}{\partial t^2} + O(k^2) .
\]
Or
\[
  \frac{\partial^2 U(x_i, 0)}{\partial t^2}
  = a^2\frac{\partial^2 U(x_i, 0)}{\partial x^2}
  = a^2\frac{\partial^2 f}{\partial x^2}
  = a^2\frac{f_{i+1} - 2f_i + f_{i-1}}{h^2} .
\]
En utilisant la condition initiale, on obtient finalement :
\begin{equation}
  U_{i,1} = (1 - \lambda^2)\,f_i + \frac{\lambda^2}{2}\bigl(f_{i+1} + f_{i-1}\bigr) + kg_i
  \tag{O6}
\end{equation}
avec $i = 1, \dots, m-1$.

Il s'agit donc de résoudre le système discret formé des équations (O2) et
(O6). La recherche des valeurs $U_{i,j+1}$ nécessite la connaissance des
valeurs $U_{i,j}$ et $U_{i,j-1}$. Au début, il faut connaître $U_{i,0}$ et
$U_{i,1}$. Les valeurs de $U_{i,0}$ découlent de (O4) et celles de $U_{i,1}$
découlent de (O6). On peut ainsi lancer le processus itératif pour calculer
toutes les valeurs de $U_{ij}$.

Il faut noter que ce schéma de discrétisation directe a une condition de
stabilité qui est
\[
  \lambda < 1 \,; \quad\text{c'est-à-dire}\quad ak/h < 1 .
\]

\begin{exercice}
\begin{enumerate}[label=\alph*-]
  \item Établir la condition de stabilité du schéma explicite ci-dessus.
  \item Établir l'algorithme du schéma explicite ci-dessus.
\end{enumerate}
\end{exercice}

\subsection{Décomposition en un système d'équations du premier ordre}

Dans certains cas, il est avantageux de décomposer l'équation d'ondes en deux
équations du premier ordre avant de discrétiser.

Par exemple, en posant
\[
  V = \frac{\partial U}{\partial x} \qquad\text{et}\qquad
  W = \frac{1}{a}\frac{\partial U}{\partial t},
\]
l'équation d'ondes devient le système d'équations suivant
\[
  \left\{
  \begin{aligned}
    \frac{\partial V}{\partial t} &= a\frac{\partial W}{\partial x}\\[0.4em]
    \frac{\partial W}{\partial t} &= a\frac{\partial V}{\partial x}
  \end{aligned}
  \right.
\]
Ce système d'équations peut être discrétisé de manière explicite et de manière
implicite.

\subsubsection{Discrétisation explicite}

Pour la discrétisation explicite, on peut utiliser la discrétisation par les
différences finies progressives dans le temps et par les différences centrées
dans l'espace. On obtient alors
\[
  \left\{
  \begin{aligned}
    V_{i,j+1} &= V_{i,j} + \frac{ak}{2h}\bigl(W_{i+1,j} - W_{i-1,j}\bigr)\\[0.3em]
    W_{i,j+1} &= W_{i,j} + \frac{ak}{2h}\bigl(V_{i+1,j} - V_{i-1,j}\bigr)
  \end{aligned}
  \right.
\]
Pour établir la condition de stabilité, posons
\[
  \left\{
  \begin{aligned}
    V_{i,j} &= A_j\,e^{Ipih}\\
    W_{i,j} &= B_j\,e^{Ipih}
  \end{aligned}
  \right.
  \qquad\text{où } I^2 = -1 .
\]
En substituant dans le système discret, on obtient
\[
  \left\{
  \begin{aligned}
    V_{i,j+1} &= \Bigl[A_j + B_j\,\frac{Iak}{h}\sin ph\Bigr]e^{Ipih}\\[0.3em]
    W_{i,j+1} &= \Bigl[B_j + A_j\,\frac{Iak}{h}\sin ph\Bigr]e^{Ipih}
  \end{aligned}
  \right.
\]
or
\[
  \left\{
  \begin{aligned}
    V_{i,j+1} &= A_{j+1}\,e^{Ipih}\\
    W_{i,j+1} &= B_{j+1}\,e^{Ipih}
  \end{aligned}
  \right.
\]
En comparant ces expressions, on obtient
\[
  \begin{pmatrix} A_{j+1}\\ B_{j+1} \end{pmatrix}
  = \begin{pmatrix} 1 & Ic\\ Ic & 1 \end{pmatrix}
    \begin{pmatrix} A_{j}\\ B_{j} \end{pmatrix}
  \qquad\text{où}\quad c = \frac{ak}{h}\sin ph .
\]
La matrice d'amplification est donc
\[
  A = \begin{pmatrix} 1 & Ic\\ Ic & 1 \end{pmatrix}.
\]
La stabilité du schéma exige que les valeurs propres de la matrice
d'amplification $A$ aient des modules inférieurs à 1. Dans le cas présent, les
modules des valeurs propres de $A$ sont tous supérieurs à 1. Donc le schéma
est inconditionnellement instable.

\subsubsection{Amélioration du schéma explicite}

Compte tenu de cette instabilité inconditionnelle, on peut améliorer le schéma
naturel ci-dessus. On peut trouver tous les $V_{i,j+1}$ de la première
équation et les utiliser ensuite dans la deuxième équation. On obtient alors le
schéma discret suivant :
\[
  \left\{
  \begin{aligned}
    V_{i,j+1} &= V_{i,j} + \frac{ak}{2h}\bigl(W_{i+1,j} - W_{i-1,j}\bigr)\\[0.3em]
    W_{i,j+1} &= W_{i,j} + \frac{ak}{2h}\bigl(V_{i+1,j+1} - V_{i-1,j+1}\bigr)
  \end{aligned}
  \right.
\]
Pour un tel schéma, la matrice d'amplification est
\[
  A = \begin{pmatrix} 1 & Ic\\ Ic & 1 - c^2 \end{pmatrix}.
\]
On prouve que le schéma est stable si et seulement si $\dfrac{ak}{h} < 2$.

\subsubsection{Schéma implicite}

On peut avoir un schéma implicite en discrétisant les dérivées spatiales à
l'instant $j+1$ et les dérivées temporelles par les différences progressives.
On obtient alors
\[
  \left\{
  \begin{aligned}
    V_{i,j+1} &= V_{i,j} + \frac{ak}{2h}\bigl(W_{i+1,j+1} - W_{i-1,j+1}\bigr)\\[0.3em]
    W_{i,j+1} &= W_{i,j} + \frac{ak}{2h}\bigl(V_{i+1,j+1} - V_{i-1,j+1}\bigr)
  \end{aligned}
  \right.
\]
Les valeurs propres de la matrice d'amplification sont dans ce cas
$1/(1 + Ic)$, $1/(1 - Ic)$.

Elles sont toutes de module $< 1$. Donc ce schéma implicite est
inconditionnellement stable.

\subsubsection{Schéma de Lax}

Si dans le schéma explicite naturel (premier schéma), on remplace $V_{i,j}$ et
$W_{i,j}$ par leurs valeurs moyennes respectives entre $i-1$ et $i+1$, alors
on obtient le schéma de Lax qui s'écrit
\[
  \left\{
  \begin{aligned}
    V_{i,j+1} &= \frac{1}{2}\bigl(V_{i+1,j} + V_{i-1,j}\bigr)
               + \frac{ak}{2h}\bigl(W_{i+1,j} - W_{i-1,j}\bigr)\\[0.3em]
    W_{i,j+1} &= \frac{1}{2}\bigl(W_{i+1,j} + W_{i-1,j}\bigr)
               + \frac{ak}{2h}\bigl(V_{i+1,j} - V_{i-1,j}\bigr)
  \end{aligned}
  \right.
\]
Dans ce cas la matrice d'amplification est
\[
  A = \begin{pmatrix} \cos ph & Ic\\ Ic & \cos ph \end{pmatrix}.
\]
La condition de stabilité est $ak/h < 1$.

\subsubsection{Schéma de Lax-Wendroff}

Pour ce schéma, on développe $V_{i,j+1}$ et $W_{i,j+1}$ jusqu'à l'ordre 2.
Par exemple
\[
  V_{i,j+1} = V_{i,j} + k\frac{\partial V_{i,j}}{\partial t}
            + \frac{k^2}{2}\frac{\partial^2 V_{i,j}}{\partial t^2}
            = V_{i,j} + ka\frac{\partial W_{i,j}}{\partial x}
            + \frac{k^2a^2}{2}\frac{\partial^2 V_{i,j}}{\partial x^2} .
\]
En faisant la même chose sur $W$, on obtient le schéma discret de
Lax-Wendroff suivant :
\[
  \left\{
  \begin{aligned}
    V_{i,j+1} &= V_{i,j} + \frac{ak}{2h}\bigl(W_{i+1,j} - W_{i-1,j}\bigr)
               + \frac{a^2k^2}{2h^2}\bigl(V_{i+1,j} - 2V_{i,j} + V_{i-1,j}\bigr)\\[0.3em]
    W_{i,j+1} &= W_{i,j} + \frac{ak}{2h}\bigl(V_{i+1,j} - V_{i-1,j}\bigr)
               + \frac{a^2k^2}{2h^2}\bigl(W_{i+1,j} - 2W_{i,j} + W_{i-1,j}\bigr)
  \end{aligned}
  \right.
\]
La matrice d'amplification est
\[
  A = \begin{pmatrix}
    1 + \dfrac{a^2k^2}{h^2}(\cos ph - 1) & \dfrac{Iak}{h}\sin ph\\[1em]
    \dfrac{Iak}{h}\sin ph & 1 + \dfrac{a^2k^2}{h^2}(\cos ph - 1)
  \end{pmatrix}.
\]
La condition de stabilité : $\dfrac{ak}{h} < 1$.

\paragraph{Remarque} Il faut noter cependant que lorsqu'une équation
hyperbolique présente des conditions initiales discontinues, la méthode des
différences finies n'est plus adaptée ; on doit faire appel à d'autres
méthodes telle que la méthode des caractéristiques.
